% Status and errata box (October 2026 freeze). \input right after \end{abstract}.
\begingroup
\setlength{\fboxsep}{7pt}%
\begin{center}
\fbox{\begin{minipage}{0.93\linewidth}\small
\setlength{\parskip}{4pt}%
\setlength{\parindent}{0pt}%
\textbf{Status and errata (October 2026).}
This is the final version of this preprint; no further revisions are planned.
It is superseded by Chapter~7 (\emph{Classification Measures}) of the monograph
\emph{Bernoulli Types: Two-Sided Error in Approximate Membership, Retrieval,
and Computation} (A.~Towell, 2026),
\url{https://metafunctor.com/latex/bernoulli-monograph/book.html}
(project hub: \url{https://metafunctor.com/research/bernoulli/}).
Please cite the monograph in preference to this preprint.

\textbf{Corrections in this version.}
\begin{itemize}\setlength{\itemsep}{1pt}\setlength{\parskip}{0pt}
\item Abstract, \cref{sec:intro}: the model is no longer called ``compositionally closed''; a two-rate pair does not propagate through arbitrary set expressions (see the companion composition paper).
\item Abstract, \cref{sec:conclusion_cm}: Bloom filters and perfect hash filters satisfy the axioms only approximately (errors share one random build); ``including Bloom filters'' is qualified.
\item \Cref{sec:prelim}: $\Expect{\beta}$ corrected from $\tprate$ to $\fnrate$ ($\beta$ is the realized false negative rate, consistent with $\mathrm{TPR}_p = 1 - \beta_p$).
\item \Cref{sec:interval_measures}: the interval PPV bounds hold for the ratio of expected counts, not for $\Expect{\PPV}$ itself.
\item \Cref{sec:intro}: novelty claims removed; the prevalence effect is the classical predictive-value relation.
\end{itemize}

\textbf{Predecessors credited.} The per-element channel is randomized
response~\cite{warner1965randomized}; misclassification in $2 \times 2$ tables is
treated by Bross~\cite{bross1954misclassification}; the dependence of predictive
values on prevalence is classical in diagnostic testing~\cite{altman1994predictive};
probabilistic precision, recall, and F-score are due to Goutte and
Gaussier~\cite{goutte2005probabilistic}; the expansions are the standard delta
method~\cite{casellaBerger}. No novelty is claimed.
A dated list of changes is in \texttt{ERRATA.md} in the source repository.
\end{minipage}}
\end{center}
\endgroup
